\section{Method}
\label{sec:method}

\begin{figure}[t]
\centering
\resizebox{\columnwidth}{!}{%
\begin{tikzpicture}[
    node distance=2.6mm and 3.2mm,
    box/.style={draw, rounded corners=2pt, align=center, minimum height=7.5mm,
                inner xsep=3pt, font=\scriptsize},
    kept/.style={box, fill=green!10},
    sub/.style={box, fill=blue!12, very thick},
    io/.style={box, fill=gray!12},
    host/.style={box, fill=orange!16},
    ar/.style={-{Latex[length=1.6mm]}, thick}
]
\node[io] (img) {image\\$1\!\times\!3\!\times\!640^2$};
\node[kept, right=of img] (bb) {preserved YOLO11\\stem, C3k2, SPPF};
\node[sub, right=of bb] (mix)
  {\textbf{S1}: C2PSA\\$\rightarrow$ C2DLA};
\node[kept, below=5.5mm of bb] (neck) {preserved PAN\\neck (P3--P5)};
\node[sub, right=of neck] (head) {\textbf{S2}: Detect\\$\rightarrow$ DetectDLA};
\node[io, right=of head] (maps) {\texttt{pred\_p3/4/5}\\$144\!\times\!\{80,40,20\}^2$};
\node[host, below=5.5mm of head] (post)
  {host task logic: sparse DFL, sigmoid, NMS, box scaling};

\draw[ar] (img) -- (bb);
\draw[ar] (bb) -- (mix);
\draw[ar] (mix.south) -- ++(0,-2.7mm) -| (neck.north);
\draw[ar] (neck) -- (head);
\draw[ar] (head) -- (maps);
\draw[ar] (maps.south) -- ++(0,-2.7mm) -| (post.north);

\node[draw=blue!65, dashed, rounded corners,
      fit=(img)(bb)(mix)(neck)(head)(maps), inner sep=3.2pt] (engine) {};
\node[font=\scriptsize\bfseries, blue!60!black, above=0.4mm of engine.north west,
      anchor=west] {DLA engine: one loadable, all tensors 4-D, static shapes};
\end{tikzpicture}%
}
\caption{Accelerator boundary and the two substitutions. Everything inside the
dashed region compiles to a single DLA loadable; everything outside it is
variable-cardinality task logic that stays on the host. S1 removes the
attention matrix multiplications, S2 removes the rank-three decoded output
interface. Green blocks are inherited from the reference model unchanged.}
\label{fig:boundary}
\end{figure}


\subsection{Execution Contract and Admissibility}
\label{sec:contract}

Let $G=(V,E)$ be a learned perception graph and let
$\mathcal{H}=(\mathcal{O},\mathcal{C})$ be a hardware execution contract:
$\mathcal{O}$ is the supported operator set and $\mathcal{C}$ collects the
configuration constraints, that is tensor rank, layout, static dimensions,
precision, kernel parameter ranges, and admissible dataflow. Write
$A_{\mathcal{H}}(G)=1$ when \emph{every} node of the learned graph is
admissible, the test being supplied by the target compiler rather than by
inspection of the source model. Offloading a detector onto a fixed-function
accelerator is then the problem of obtaining an admissible graph that stays as
close as possible to a reference $G_0$ with function $f_0$, which we treat as a
Pareto problem over structural change and behavioral deviation,
\begin{equation}
\min_{G_R}\ \left(d_{\mathrm{arch}}(G_0,G_R),\,
D_{\mathrm{ref}}(f_0,f_R)\right)
\quad\text{s.t.}\quad A_{\mathcal{H}}(G_R)=1,
\label{eq:retrofit}
\end{equation}
where $d_{\mathrm{arch}}$ measures structural change and $D_{\mathrm{ref}}$
measures deviation under the task metrics and, where required, output agreement
on representative inputs. Equation~\eqref{eq:retrofit} is a design-selection
criterion, not an additional term in the detector's training loss. Its
constraint is binary and supplied by the toolchain, which is what makes the
formulation transferable: on an NPU, an Edge TPU, or an FPGA toolchain the same
criterion applies with a different $\mathcal{H}$ and a different acceptance
procedure. Here $G_0$ is YOLO11 \cite{jocher2024yolo11} and $\mathcal{H}$ is the
NVIDIA DLA contract \cite{nvidia_dla_restrictions,nvidia_dla_build}, whose
relevant clauses are that tensors be four-dimensional with static shapes, that
no matrix multiplication appear, that slices and concatenations act on rank-four
tensors, and that broadcasts be NCHW-compatible.

\subsection{The Accelerator Boundary}
\label{sec:boundary}

Admissibility is not only a question of which operators appear; it is a question
of where the learned graph is declared to end. A detector's terminal
post-processing is inherently variable in cardinality: the number of surviving
candidates depends on image content, and the standard formulation reaches it by
flattening the pyramid into a rank-three tensor of $8400$ anchors, decoding
distributions, building an anchor grid, and suppressing overlaps. None of that
is admissible under $\mathcal{H}$, and worse, its inadmissibility is
contagious. A compiler that partitions a graph by fusible regions will pull
otherwise admissible producers into the fallback partition that feeds its
rank-changing consumer, which is precisely the failure we characterize in
Sec.~\ref{sec:rq2}.

We therefore fix the boundary at static, raw, multi-scale predictions. The
compiled graph accepts a fixed tensor
$I\in\mathbb{R}^{B\times3\times H\times W}$, keeps every intermediate tensor
four-dimensional, and for each pyramid level $\ell\in\{3,4,5\}$ returns
\begin{equation}
P_\ell=\operatorname{Concat}_{C}(R_\ell,S_\ell)
\in\mathbb{R}^{B\times(4r+n_c)\times H_\ell\times W_\ell},
\label{eq:packed}
\end{equation}
where $R_\ell$ holds four $r$-channel distance predictions, $S_\ell$ holds
$n_c$ class logits, and the concatenation is over channels, which the contract
allows on rank-four tensors. No flattening, cross-scale concatenation, anchor
construction, top-$k$, or non-maximum suppression (NMS) is part of the compiled
graph. The host then restores the scale-wise grids and decodes each side
$q\in\{l,t,r,b\}$ by the distributional expectation
\begin{equation}
\hat d_q=\sum_{j=0}^{r-1}j\,
\frac{\exp(R_{qj})}{\sum_{k=0}^{r-1}\exp(R_{qk})},
\label{eq:dfl}
\end{equation}
converts distances relative to cell centers, multiplies by the pyramid stride,
takes $\sigma(S_\ell)$ as class probabilities, and applies NMS.

This boundary is the load-bearing design decision of the paper, and it has three
consequences that we measure separately. First, it is what makes a single
fallback-free loadable possible at all: with the decode inside the graph, no
amount of operator substitution upstream produces an admissible model. Second,
it moves work off the accelerator, so an engine-boundary comparison against a
conventional detector is not a comparison of equal pipelines, since the
reference model emits one decoded $1\times84\times8400$ tensor while the
retrofit emits three raw maps and pays for their decode on the host; we
therefore report a stage-resolved accounting (Sec.~\ref{sec:rq3}) rather than
engine latency alone. Third, it exposes an interface in an accelerator-native
layout: a fixed-function engine is free to demand a packed input format and to
return channel-vectorized outputs, and the cost of meeting that demand is paid
outside the engine, by the host, in whatever implementation the integrator
supplies. That cost is real, it is large in our reference implementation, and it
is not a property of the accelerator; Sec.~\ref{sec:packing} characterizes it
explicitly.

\subsection{YOLO11-DLA-n}
\label{sec:retrofit}

\begin{figure}[t]
\centering
\resizebox{\columnwidth}{!}{%
\begin{tikzpicture}[
    op/.style={draw, rounded corners=2pt, align=center, minimum height=6.4mm,
               minimum width=8.5mm, inner xsep=2.5pt, font=\scriptsize},
    conv/.style={op, fill=blue!10},
    act/.style={op, fill=violet!12},
    io/.style={op, fill=gray!14, minimum width=6mm},
    sum/.style={draw, circle, inner sep=0.5pt, font=\scriptsize},
    ar/.style={-{Latex[length=1.6mm]}, thick}
]
\node[io]   at (0.25,-0.65) (b)   {$X$};
\node[conv] at (1.45,0)     (pk)  {$\phi_k$};
\node[act]  at (2.75,0)     (sm)  {softmax\\over $C$};
\node[conv] at (4.05,0)     (pa)  {$\phi_a$};
\node[conv] at (1.45,-1.30) (pv)  {$\phi_v$};
\node[conv] at (2.75,-1.30) (dw)  {DW $3\!\times\!3$};
\node[sum]  at (4.95,-0.65) (s1)  {$+$};
\node[conv] at (5.75,-0.65) (po)  {$\phi_o$};
\node[sum]  at (6.55,-0.65) (r1)  {$+$};
\node[conv] at (7.40,-0.65) (ffn) {FFN};
\node[sum]  at (8.25,-0.65) (r2)  {$+$};
\node[io]   at (9.00,-0.65) (out) {$X''$};

\draw[ar] (b.east) -- ++(0.18,0) |- (pk.west);
\draw[ar] (b.east) -- ++(0.18,0) |- (pv.west);
\draw[ar] (pk) -- (sm);
\draw[ar] (sm) -- (pa);
\draw[ar] (pv) -- (dw);
\draw[ar] (pa.east) -| (s1.north);
\draw[ar] (dw.east) -| (s1.south);
\draw[ar] (s1) -- (po);
\draw[ar] (po) -- (r1);
\draw[ar] (r1) -- (ffn);
\draw[ar] (ffn) -- (r2);
\draw[ar] (r2) -- (out);
% residual around the mixer, Eq. (8)
\draw[ar] (0.25,-0.34) -- (0.25,0.80) -- (6.55,0.80) -- (r1.north);
% residual around the FFN, Eq. (9)
\draw[ar] (6.55,-0.96) -- (6.55,-1.95) -- (8.25,-1.95) -- (r2.south);
\end{tikzpicture}%
}
\caption{One C2DLA mixing block, i.e.\ Eqs.~\eqref{eq:kv}--\eqref{eq:res2}.
Only this branch differs from the reference C2PSA block: the pairwise spatial
attention becomes a rank-$d$ channel-softmax bottleneck built from $1\times1$
convolutions, while the positional depthwise convolution on $V$, both
residuals, and the feed-forward block (two $1\times1$ convolutions with SiLU
between them) are inherited. $X$ is the processed half of the enclosing
split--process--merge structure, which is also inherited; independent $\phi_k$
and $\phi_v$ replace the reference fused QKV projection so that no channel
slice and no shape-helper subgraph reaches the compiler.}
\label{fig:c2dla}
\end{figure}


\subsubsection{Two substitutions}
Under $\mathcal{H}$, exactly two regions of YOLO11n are inadmissible: the
position-sensitive attention in the terminal backbone stage, which forms a
spatial affinity matrix by matrix multiplication, and the decoded output
interface discussed above. The retrofit is therefore
\begin{equation}
\mathrm{C2PSA}\rightarrow\mathrm{C2DLA},\qquad
\mathrm{Detect}\rightarrow\mathrm{DetectDLA},
\label{eq:substitutions}
\end{equation}
and nothing else changes: the stem, all C3k2 stages, SPPF, the full
path-aggregation neck, the channel schedule, the P3--P5 connections, the head
towers, task-aligned one-to-many assignment, the losses, and the $r=16$
distributional regression \cite{li2020gfl} are inherited (Fig.~\ref{fig:boundary},
Table~\ref{tab:design}). What we compare against the reference is thus a
structural approximation of it, not a new detector.

\subsubsection{C2DLA}
C2DLA preserves C2PSA's split--process--merge structure and its residuals, and
replaces only the attention math (Fig.~\ref{fig:c2dla}). Let
$X\in\mathbb{R}^{B\times C\times H\times W}$ be the processed half. One block
computes
\begin{align}
K &= \phi_k(X), \qquad V = \phi_v(X), \label{eq:kv}\\
A_{bkhw} &= \frac{\exp(K_{bkhw})}{\sum_{j=1}^{d}\exp(K_{bjhw})},
  \label{eq:chansoftmax}\\
Z &= \phi_o\!\left(\phi_a(A)+
  \operatorname{DWConv}_{3\times3}(V)\right), \label{eq:mix}\\
X' &= X+Z, \label{eq:res1}\\
X'' &= X'+\phi_2\!\left(\operatorname{SiLU}(\phi_1(X'))\right),
  \label{eq:res2}
\end{align}
with $\phi_k:C\rightarrow d$, $\phi_v:C\rightarrow C$, $\phi_a:d\rightarrow C$
and $\phi_o:C\rightarrow C$ all $1\times1$ convolutions and
$d=\max(C/4,8)$; the unprocessed half and $X''$ are concatenated and projected
as in C2PSA. Three properties matter. The normalization in
\eqref{eq:chansoftmax} runs over channels at each spatial location, so no
affinity matrix and no transpose is formed, and the block follows external
attention in replacing pairwise interaction by projections through a compact
memory of rank $d$ \cite{guo2023external}. The mixing in \eqref{eq:mix} keeps
the reference block's positional depthwise convolution on $V$, which is the
component that supplies local spatial context. Finally, $\phi_k$ and $\phi_v$
are separate convolutions rather than a fused QKV projection followed by a
channel slice: the fused form is equivalent in function but exports as a
\texttt{Slice} chain surrounded by shape helpers, which the compiler places on
the GPU and which then costs a round trip off the accelerator per scale. The
price of the substitution is that mixing is no longer data-dependent or
multi-head; it is a learned low-rank channel mixer gated by a channel softmax.

\subsubsection{DetectDLA}
DetectDLA leaves training untouched: it uses the original box and class towers
and the standard rank-three path in training and in ordinary PyTorch inference.
Only in export mode does it bypass DFL, anchor decoding, sigmoid, and NMS and
emit one packed tensor per scale according to \eqref{eq:packed}. For COCO
($n_c=80$, $r=16$) at $640\times640$ the three engine outputs are
$1\times144\times80\times80$, $1\times144\times40\times40$, and
$1\times144\times20\times20$, exposed as the bindings \texttt{pred\_p3},
\texttt{pred\_p4}, \texttt{pred\_p5} ordered by stride. Export metadata records
the number of raw outputs, the strides, the regression width $r$, the class
names, and whether the head is one-to-one, so that the host decoder is selected
from the engine's declared contract instead of being hard-coded.

\subsubsection{Weight-aware initialization}
Because the edit is localized, most of the reference model transfers directly.
All shape-compatible parameters outside the substituted block are loaded from a
task-matched YOLO11 checkpoint. Inside the block, the positional depthwise
convolution and the output projection are copied verbatim; the per-head
channel slices of the reference fused QKV projection corresponding to $K$ and
$V$, together with their batch-normalization statistics, are gathered into
$\phi_k$ and $\phi_v$ whenever the shapes agree, which for $\phi_k$ requires
$d$ to equal the reference key width; and $\phi_a$, which has no counterpart in
the reference, is initialized randomly. Training then proceeds in two phases:
the layers preceding the substituted block are frozen while the new projection
adapts, after which the whole model is fine-tuned at a lower learning rate.
Initialization is part of the method and not an implementation detail, since it
is what allows a substituted block to approximate a reference block's behavior
rather than merely occupy its position.

\begin{table}[t]
\caption{Reference detector and its admissible retrofit. Only the two rows in
bold differ by construction; every other row is inherited.}
\label{tab:design}
\centering
\scriptsize
\setlength{\tabcolsep}{2.5pt}
\renewcommand{\arraystretch}{1.1}
\begin{tabular}{p{0.20\columnwidth}p{0.34\columnwidth}p{0.36\columnwidth}}
\toprule
Axis & YOLO11n (reference) & YOLO11-DLA-n (retrofit) \\
\midrule
Stem, C3k2 stages, SPPF, PAN neck & inherited & identical \\
\textbf{Terminal backbone stage} & C2PSA, position-sensitive attention with
two matrix multiplications & C2DLA, rank-$d$ channel-softmax mixer, $1\times1$
convolutions only \\
Positional encoding & $3\times3$ depthwise on $V$ & unchanged \\
Activations & SiLU & unchanged \\
Box representation & $r=16$ distributional & unchanged \\
Assignment, losses & task-aligned one-to-many & unchanged \\
\textbf{Exported interface} & one decoded $1\times84\times8400$ tensor &
three raw maps, $4r+n_c=144$ channels per scale \\
Host task logic & NMS & DFL decode, sigmoid, NMS \\
Initialization & official COCO checkpoint & shape-compatible transfer, block
hot-start, two-phase fine-tune \\
Strict DLA build & rejected at the attention \texttt{MatMul} & accepted, one
loadable, no layer on GPU \\
\bottomrule
\end{tabular}
\end{table}
